\documentclass[12pt,a4paper,fleqn]{article}

%%%%%%%%%%%%%%%%%%%%%% LES PACKAGES
%-------------------
% gestion des pages
%-------------------
\usepackage[left=1.5cm,right=1.5cm,top=1cm,bottom=2cm]{geometry}
\usepackage{pdflscape} %pages horizontales dans le PDF
\usepackage{calc} %pour pouvoir écrire \textwidth-1cm

%--pour les numéro de pages/ pieds de pages
%-- voir mise_en_forme_MP2I.tex pour mise en oeuvre
\usepackage{fancyhdr}


%-------------------
% sommaire 
%-------------------
\usepackage{tocloft}%mise en forme de la toc
\usepackage{titlesec}%mise en forme des titres de section
\usepackage[pagebackref,hidelinks]{hyperref}%pour les liens hypertextes


%-------------------
% langue et encodage
%-------------------
\usepackage[utf8]{inputenc}
\usepackage[T1]{fontenc}
\usepackage[]{babel}
%\usepackage[english]{babel}
%\usepackage[french]{babel}


%-------------------
% symboles et mode maths
%-------------------
\usepackage{amsmath,amsfonts,amssymb}
\usepackage{textcomp,lmodern}%pour l'euro
\usepackage{mathrsfs}%lettres style manuscrit italique


%-------------------
% code
%-------------------
\usepackage{minted}
%%-- pour la commande qui évite l'italique
%%-- voir https://github.com/gpoore/minted/issues/71
%\usepackage{etoolbox}
%\usepackage{newunicodechar}



%-------------------
% environnement, théorèmes et exercices
%-------------------
%\usepackage{thmbox}%pour les jolis thm


%-------------------
% tableaux
%-------------------
\usepackage{array,multirow}


%-------------------
% images
%-------------------
\usepackage{xcolor} %pour redefinir des couleurs
\usepackage{caption}%les personaliser (les centrer par ex)
\usepackage{graphicx} % pour includegraphics

\usepackage{tikz} %graphiques et dessins
\usetikzlibrary{shapes}%pr écrire \node[ellipse] par ex
\usetikzlibrary{positioning} %pr écrire \node[below rigth = 3pt and 5pt]
\usetikzlibrary{decorations.pathreplacing}%pr les accolades
\usetikzlibrary{patterns}%pour les hachures





%-------------------
% divers
%-------------------
\usepackage{cancel}%pour barrer
%-------------------
\usepackage{lipsum}% for dummy text
%-------------------



\usepackage[most]{tcolorbox}
\usepackage[version=4]{mhchem}

%%%%%%%%%%%%%%%%%%%%%% MISE EN FORME
%%%%%%%%%%%%%%%%%%%%%%%%%%%%%%%%%%%%%%%%%%%%
% MISE EN FORME DU TITRE 
%%%%%%%%%%%%%%%%%%%%%%%%%%%%%%%%%%%%%%%%%%%%
\newcommand{\titreCopie}[1]{
\noindent
%\\%[0.4cm]
\begin{center}
\LARGE
\hrule
\vspace{.4cm}
#1 \\[0.4cm]
\hrule
\end{center}

\vspace*{0.4cm}
}


%%%%%%%%%%%%%%%%%%%%%%%%%%%%%%%%%%%%%%%%%%%%
% MISE EN FORME DES PIEDS DE PAGE
%%%%%%%%%%%%%%%%%%%%%%%%%%%%%%%%%%%%%%%%%%%%
\AtEndDocument{\label{lastpage}}
\lhead{}
\chead{}
\rhead{}
\lfoot{ 
Ryan Bouchou
}
\cfoot{}
\rfoot{\thepage/\pageref{lastpage}}
\renewcommand{\headrulewidth}{0pt}
\fancyhfoffset{0.7cm}



%%%%%%%%%%%%%%%%%%%%%%%%%%%%%%%%%%%%%%%%%%%%
% REMISE EN FORME DES SOUS-SOUS-SECTIONS  
%%%%%%%%%%%%%%%%%%%%%%%%%%%%%%%%%%%%%%%%%%%%

% section = exercices
\renewcommand{\thesection}{Exercice~\arabic{section}}
% sous-section = question
\renewcommand{\thesubsection}{Question~\arabic{subsection}}




%%%%%%%%%%%%%%%%%%%%%%%%%%%%%%%%%%%%%%%%%%%%
% MISE EN FORME DU CODE AVEC MINTED
%%%%%%%%%%%%%%%%%%%%%%%%%%%%%%%%%%%%%%%%%%%%

%pour le style des numéros de ligne minted
\renewcommand{\theFancyVerbLine}{\sffamily
\textcolor{expli}{\scriptsize
\oldstylenums{\arabic{FancyVerbLine}}}}

%pour éviter l'italique 
%car apparaissaient en italique à la fois les commentaires et les includes
\newtoggle{inminted}
\AtBeginEnvironment{minted}{\let\itshape\relax}

%paramètres pour le c
\setminted[c]{
	%--fond et cadre	
	%bgcolor = black!3!white,
	frame = leftline,
	framesep = 6pt,
	rulecolor= expli,
	%--numéro de ligne
	linenos=true,
	numbersep=4pt,
	%stepnumber=2,
	xleftmargin=20pt,%pour que les chiffres débordent pas
	%--découpage des longues lignes
	breaklines=true,
	%--tabulations
	tabsize=2,
	}

%paramètres pour le ocaml
\setminted[ocaml]{
	%--fond et cadre	
	%bgcolor = black!3!white,
	frame = leftline,
	framesep = 6pt,
	rulecolor= expli,
	%--numéro de ligne
	linenos=true,
	numbersep=4pt,
	%stepnumber=2,
	xleftmargin=20pt,%pour que les chiffres débordent pas
	%--découpage des longues lignes
	breaklines=true,
	%--tabulations
	tabsize=2,
	}




%%%%%%%%%%%%%%%%%%%%%%% LA PALETTE DE COULEURS
%%%%%%%%%%%%%% couleurs locales
\definecolor{jaune}{RGB}{255, 195, 0 }
\colorlet{oranger}{orange!50!jaune}
\definecolor{orange}{RGB}{250, 130, 1}

\definecolor{vert}{RGB}{194, 247, 50}
\definecolor{herbe}{RGB}{80, 180, 33}
\definecolor{vertFonce}{RGB}{20, 148, 20} 

\colorlet{vertdEau}{blue!50!green!30!white} 
\definecolor{turquoise}{RGB}{0, 204, 203}
\colorlet{monCyan}{turquoise!90!blue} 
\definecolor{turquoiseFonce}{RGB}{0, 149, 182} 
\colorlet{marine}{blue!60!black}
\colorlet{doubleu}{gray!50!blue}


\colorlet{lavande}{vertdEau!50!blue}
\definecolor{mauve}{RGB}{161, 132, 220}
\definecolor{lilas}{RGB}{154, 107, 165}


\definecolor{prunelle}{RGB}{106,24, 141}
\definecolor{prune}{RGB}{121, 7, 123}
\definecolor{violet}{RGB}{153, 0, 139}
\definecolor{violette}{RGB}{128, 0, 128} 

\definecolor{roseAcidule}{RGB}{255, 0, 127} 
\colorlet{rose}{orange!50!oranger!50!magenta!50!violet}
\definecolor{grenat}{RGB}{141, 24, 59}
\colorlet{framboise}{magenta!90!red!80!black}




\colorlet{ocre}{orange!90!green!80!white} 
\definecolor{chamois}{RGB}{200, 141, 75}

\definecolor{briqueRouge}{RGB}{200, 48, 24}
\definecolor{brique}{RGB}{141, 48, 24}
\definecolor{taupe}{RGB}{165, 137, 107}


\colorlet{expli}{gray}


\colorlet{saumon}{yellow!45!magenta}



%%%%%%%%%%%%%%%%%%%%%%% LES COMMANDES 
%%%%%%%%%%%%%%%%%%%%%%%%%%%%%%%%%%%%%%%%%%%%
% GENERAL  
%%%%%%%%%%%%%%%%%%%%%%%%%%%%%%%%%%%%%%%%%%%%

%raccourcis généraux
\newcommand {\ie}{\textit{i.e}. }
\newcommand {\cf}{\textit{Cf.} }
\newcommand {\resp}{resp.\! }


%pour les belles fonctions (à mettre entre $ $  ou entre $$ $$ )
\newcommand{\fonction}[5]{ 
#1 =
\left( \! \begin{tabular}{c@{ }c@{ }l} 
$#2$  & $\longrightarrow$   & $#3$  \\
$#4$  & $\longmapsto$       & $#5$ 
\end{tabular} \! \right)
}

%le widebar (inutile si package mathabx
\newcommand{\widebar}[1]{\overline{#1}}

%symboles rapprochés
\renewcommand {\=}{\!=\!}
\renewcommand {\-}{\!-\!}
\renewcommand {\:}{\!:\!}
\newcommand {\+}{\!+\!}
\newcommand {\iin}{\!\in\!}
\newcommand {\inclus}{\!\subseteq\!}

%symboles
\newcommand{\bubullet}{-}
\newcommand{\subbullet}{\circ}
\newcommand{\parties}{\mathcal{P}}

%produit scalaire
\renewcommand{\.}{\!\cdot\!}

%indicatrice
\newcommand{\1}{\mathbb{I}}

%cardinal
\newcommand{\card}[1]{\left|#1\right|}

%%%%%%%%%%%%%%%%%%%%%%%%%%%%%%%%%%%%%%%%%%%%
% ENSEMBLES DE NOMBRES & ESPACES DE VARIABLES  
%%%%%%%%%%%%%%%%%%%%%%%%%%%%%%%%%%%%%%%%%%%%

%pour N,Z,R 
\newcommand {\NN}{\mathbb{N}}
\newcommand {\NNe}{\NN^*}
\newcommand {\R}{\mathbb{R}}
\newcommand {\Rp}{\R_+}
\newcommand {\Rpe}{\R^*_+}
\renewcommand {\Re}{\R^*}
\newcommand {\Z}{\mathbb{Z}}


%%%%%%%%%%%%%%%%%%%%%%%%%%%%%%%%%%%%%%%%%%%%
% LES NOUVEAUX OPERATEURS  
%%%%%%%%%%%%%%%%%%%%%%%%%%%%%%%%%%%%%%%%%%%%

%optim
\DeclareMathOperator*{\argmin}{arg\,min}
\DeclareMathOperator*{\argmax}{arg\,max}

%analyse cvx
\DeclareMathOperator*{\conv}{conv}
\DeclareMathOperator*{\extr}{extr}
\DeclareMathOperator*{\vect}{vect}
\DeclareMathOperator*{\aff}{aff}





%%%%%%%%%%%%%%%%%%%%%%%%%%%%%%%%%%%%%%%%%%%%
% ICONES ET SYMBOLES HORS MATHS 
%%%%%%%%%%%%%%%%%%%%%%%%%%%%%%%%%%%%%%%%%%%%
\newcommand{\flch}{\item[$\rightarrow$]}

\newcommand{\attention}{\dbend}

\newcommand{\alamain}{{\LARGE\ding{45}}}

\newcommand{\surmachine}{
\begin{tikzpicture}[scale=0.1]
\draw[line width=1.2pt](-2,0) rectangle (2,3);
\fill(-1,0)--(1,0)--(2,-1)--(-2,-1)--(-1,0);
\end{tikzpicture}
}

\newcommand{\pasapas}{
\begin{tikzpicture}[scale=0.1]
\draw (0,0)--(0,1)--(1,1)--(1,2)--(2,2)--(2,3)--(3,3)--(3,4)--(4,4)--(5,4)--(5,0)--(0,0);
\end{tikzpicture}
}

\newcommand{\afaire}{
\begin{tikzpicture}[scale=0.1]
\draw (0,0)rectangle(4,4);
\end{tikzpicture}
}

\newcommand{\uneetoile}{{\LARGE$*$}\,}
\newcommand{\deuxetoiles}{{\LARGE$*\,*$}\,}
\newcommand{\troisetoiles}{{\LARGE$_*^*$}{\large$*$}\,}

\newtcolorbox{coder}{
enhanced,
boxrule=0pt,frame hidden,
borderline west={4pt}{0pt}{briqueRouge},
colback=black!3!white,
sharp corners
}
\newtcolorbox{codeb}{
enhanced,
boxrule=0pt,frame hidden,
borderline west={4pt}{0pt}{vertdEau},
colback=black!3!white,
sharp corners
}

\PassOptionsToPackage{svgnames}{xcolor}

%%%%%%%%%%%%%%%%%%%%%%%%%%%%%%%%%%%%%%%
\begin{document}
\pagestyle{fancy} %active les pieds de pages
\tcbset{colback=red!5!white,colframe=red!75!black,fonttitle=\bfseries}


%%%%%%%%%%%%%%%%%%%%%%%%%%%%%%%%%%%%%%%
\titreCopie{TD Maths}
%%%%%%%%%%%%%%%%%%%%%%%%%%%%%%%%%%%%%%%
\section{Calcul littéral}
\subsection{Simplification d'expressions}
\begin{enumerate}
    \item \(3x + 2y - 5x + 4y\)
    \item \(2a - 3b + 5a + b\)
    \item \(4m - 2n + 7m - 3n\)objet
\end{enumerate}
\subsection{Factorisation}
\begin{enumerate}
    \item  \(6x + 9y\)
    \item  \(15a - 5b\)
    \item  \(10m^2 - 5n\)
\end{enumerate}
\subsection{Développement}
\begin{enumerate}
    \item \((x + 3)(x - 2)\)
    \item \((2a - 5)(3a + 4)\)
    \item \((3m + 2)(m - 1)\)
\end{enumerate}
\subsection{Résolution d'équations}
\begin{enumerate}
    \item \(2x + 5 = 3x - 1\)
    \item \(4(x - 2) = 12\)
    \item \(2(3y + 1) = 5y - 3\)
\end{enumerate}
\subsection{Problème}
Un rectangle a une aire $A$ donnée par \(A = 3x^2 + 5x - 2\). Si la longueur du rectangle est \(x + 2\), quelle est l'expression pour la largeur du rectangle ? Quelle est la largeur si $x=3$ ?

\section{Dérivation}
\subsection{Dérivées de fonctions simples}
\begin{enumerate}
 \item \(f(x) = 3x^2 - 4x + 1\)
 \item \(g(x) = \sqrt{x} + 2x\)
 \item \(h(x) = e^x - \ln(x)\)
\end{enumerate}
\subsection{Règles de dérivation}
\begin{enumerate}
    \item  \(f(x) = 4x^3 + 2x^2 - 5x + 1\)
    \item  \(g(x) = \frac{2}{x} + \sin(x)\)
    \item  \(h(x) = e^{2x} \cdot \cos(x)\)
\end{enumerate}
\subsection{Dérivées successives}
\begin{enumerate}
    \item \(f(x) = x^4 - 2x^3 + 3x^2 - x + 1\)
    \item  \(g(x) = \frac{1}{x^2} + \cos(x)\)
\end{enumerate}
\subsection{Problème}
\subsubsection{} La fonction de coût d'une entreprise est donnée par \(C(x) = 2x^2 + 3x + 5\), où \(x\) est la quantité produite. Trouver la dérivée de la fonction coût par rapport à \(x\).
\subsubsection{} La position d'une particule est donnée par la fonction \(s(t) = 4t^2 - 3t + 1\), où \(t\) est le temps en secondes. Trouver la vitesse de la particule en fonction du temps en calculant la dérivée de \(s(t)\).
\subsection{}
Soit \(f(x) = x^3 - 3x^2 + 2x + 5\). Trouver les points où la tangente à la courbe de \(f(x)\) est horizontale.
\section{Fonction $e^x$ et $ln$}
\subsection{Exponentielle}
\begin{enumerate} 
\item  Soit \(f(x) = 2^x\). Calculer \(f(0)\), \(f(1)\), et \(f(2)\).
\item  Résoudre l'équation \(2^x = 8\).
\item  Simplifier l'expression \(\frac{2^{x+1}}{2^x}\).
\end{enumerate}
\subsection{Logarithme}
\begin{enumerate}
\item Résoudre l'équation \(\log_2(x) = 3\).
\item Simplifier l'expression \(\log_3(9) - \log_3(3)\).
\item Si \(\log_a(b) = 2\), quelle est la valeur de \(a\)?
\end{enumerate}
\subsection{Applications}
\subsubsection{} 
La population d'une ville suit une croissance exponentielle de \(P(t) = 1000 \times 1.02^t\), où \(t\) est le nombre d'années depuis aujourd'hui. Quelle sera la population dans 5 ans?
\subsubsection{} Un investissement initial de 5000 € augmente de 8\% chaque année. Modélisez la valeur de l'investissement en fonction du temps avec une fonction exponentielle.
\subsubsection{} Une substance radioactive perd 10\% de son activité chaque année. Si l'activité initiale est de 1000 Bq, quelle sera l'activité après 3 ans?
\subsection{Plusieurs fonctions}
\begin{enumerate}
    \item  Soit \(g(x) = e^{2x}\) et \(h(x) = \ln(x)\). Calculer \(g(1)\) et \(h(e)\).
    \item  Résoudre l'équation \(e^{2x} = 16\).
\end{enumerate}
\subsection{}
\begin{enumerate}
    \item  Représenter graphiquement la fonction \(y = e^x\) sur l'intervalle \(\mathbb{R}\) ainsi que les éventuelles asymptotes, et la tangeante en 0.
\item  Déterminer la ou les asymptotes de la fonction \(y = \ln(x)\).
\item  Soit \(k(x) = \frac{e^{2x} - 1}{e^x + 1}\). Simplifier \(k(x)\).
$$log_b(x)=\frac{ln(x)}{ln(b)}$$
\end{enumerate}
\end{document}

